%% ma: L10-B14-0073
%% lop: 10
%% chuong: 5
%% bai: 14
%% dang: 10.14.8
%% loai: TN4
%% muc: VD
%% dap_an: "A"
%% nguon: "chuyen-de-cac-so-dac-trung-cua-mau-so-lieu-khong-ghep-nhom-toan-10-chuong-trinh-moi.pdf (GV. Phan Nhật Linh) (P05), Bài 03 Dạng 2 Phần I Câu 18"
%% kien_thuc: "Lớp 10 Bài 13 (số trung bình), Bài 14 (phương sai, độ lệch chuẩn)"
%% trang_thai: chinh-thuc
%% ghi_chu: "Tự duyệt theo yêu cầu giáo viên 22/10/2026."
\begin{ex}
Bảng sau cho biết thời gian làm bài (phút) của $50$ học sinh:
\begin{center}
\begin{tabular}{|c|c|c|c|c|c|c|c|c|c|c|}\hline
Thời gian & $3$ & $4$ & $5$ & $6$ & $7$ & $8$ & $9$ & $10$ & $11$ & $12$\\\hline
Tần số & $1$ & $3$ & $4$ & $7$ & $8$ & $9$ & $8$ & $5$ & $3$ & $2$\\\hline
\end{tabular}
\end{center}
Độ lệch chuẩn của mẫu số liệu (làm tròn đến hàng phần trăm) là
\choice{\True $2{,}13$}{$2{,}14$}{$2{,}16$}{$2{,}15$}
\loigiai{$\overline x=\dfrac{384}{50}=7{,}68$; $\dfrac1{50}\sum n_ix_i^2=63{,}52$.

$s^2=63{,}52-7{,}68^2=4{,}5376$; $s\approx2{,}13$. Chọn A.}
\end{ex}
